% Chapter - Stochasticity and Markov Processes ................................................................

\chapter{Stochasticity and Markov Processes}
\label{chapter7}

\epigraph{\textit{The theory of probabilities is at bottom nothing\\but common sense reduced to calculation.}}{— Pierre-Simon Laplace}

The mathematization of uncertainty began with the analysis of games of chance, where Bernoulli’s law of large numbers established that stable regularities emerge from repeated random trials, revealing that randomness itself could obey precise quantitative laws and need not imply disorder \cite{bernoulli_1713_ars}. 

\medskip

Laplace extended this insight beyond gambling to scientific reasoning, arguing that probability expresses rational inference under incomplete knowledge and that apparent randomness often reflects limited information rather than metaphysical indeterminacy, thereby embedding uncertainty within a broader epistemological framework \cite{laplace_1812_theorie}. 

\medskip

During the nineteenth and early twentieth centuries, attention shifted from isolated probabilistic events to sequences of dependent variables evolving over time, culminating in Markov’s demonstration that dependence need not destroy probabilistic regularity, as he extended the law of large numbers to dependent trials and initiated the systematic study of stochastic processes \cite{markov_1906_extension}. 

\medskip

This temporal generalization was later placed on rigorous foundations through Kolmogorov’s axiomatization of probability and Doob’s development of stochastic process theory, which formalized random evolution as measurable structure on a probability space and established the analytical framework that continues to underlie modern stochastic analysis \cite{kolmogorov_1933_grundbegriffe,doob_1953_stochastic}. 

\medskip

Scientific models often begin with deterministic equations, yet many natural systems exhibit variability that cannot be captured by fixed trajectories alone, and stochastic modelling therefore represents not chaos but structured uncertainty governed by probabilistic laws, clarifying that variability over time or across states may be intrinsic rather than merely observational and requiring explicit probabilistic structure rather than deterministic approximation. 

% Section .....................................................................................................

\section{Mathematical definition}

Let \((\Omega,\mathcal{F},\mathbb{P})\) be a probability space; a stochastic process is a family of random variables 
\[
\{X_t : t \in T\},
\]
defined on this common probability space, where the index set \(T\) typically represents time. 

\medskip

Each \(X_t\) is a measurable function \(X_t : \Omega \to S\), where \(S\) is the state space, and the joint distribution of finite collections \((X_{t_1},\dots,X_{t_n})\) determines the probabilistic behavior of the process \cite{kolmogorov_1933_grundbegriffe}. 

\medskip

Dependence structure distinguishes stochastic processes from independent sequences; in general,
\[
\mathbb{P}(X_{t+1} \in A \mid X_t, X_{t-1},\dots) 
\]
may depend on the entire past history. 

\medskip

The Markov property simplifies this dependence by asserting that
\[
\mathbb{P}(X_{t+1} \in A \mid X_t, X_{t-1},\dots) 
=
\mathbb{P}(X_{t+1} \in A \mid X_t),
\]
meaning that the future is conditionally independent of the past given the present state. 

\medskip

This conditional independence assumption captures systems with “memoryless” dynamics, reducing complex temporal dependence to tractable transition structure \cite{doob_1953_stochastic}. 

% Section .....................................................................................................

\section{Stochasticity and Markov processes}

A Markov process is therefore a stochastic process satisfying the Markov property, and its behavior is characterized by transition probabilities between states. 

\medskip

In discrete time and finite state space \(S = \{1,\dots,k\}\), these transitions are described by a matrix \(P = (p_{ij})\) with entries
\[
p_{ij} = \mathbb{P}(X_{t+1}=j \mid X_t=i),
\]
where each row sums to one. 

\medskip

The distribution of the process evolves according to
\[
\pi_{t+1} = \pi_t P,
\]
where \(\pi_t\) is the row vector of state probabilities at time \(t\), linking stochastic dynamics directly to linear algebraic structure. 

\medskip

Under suitable conditions such as irreducibility and aperiodicity, a Markov chain possesses a stationary distribution \(\pi\) satisfying 
\[
\pi = \pi P,
\]
representing long-run equilibrium behavior independent of initial state. 

\medskip

Markov processes thus provide a mathematically precise model of evolving uncertainty, capturing dependence while retaining analytical tractability. 

% Section .....................................................................................................

\section{Markov chains and Monte Carlo methods}

Markov chains play a central role in modern computational statistics through Markov Chain Monte Carlo (MCMC) methods, which generate dependent samples from complex probability distributions. 

\medskip

Suppose a target distribution \(\pi(\theta)\) is difficult to sample directly; MCMC constructs a Markov chain with transition kernel \(K(\theta,\theta')\) such that \(\pi\) is stationary, meaning 
\[
\int \pi(\theta)K(\theta,\theta')\,d\theta = \pi(\theta').
\]

\medskip

Under mild regularity conditions, ergodic theorems ensure that sample averages along the chain converge to expectations under \(\pi\), extending classical law of large numbers arguments to dependent sequences \cite{markov_1906_extension}. 

\medskip

Algorithms such as the Metropolis–Hastings method achieve this by proposing candidate moves and accepting them with probability
\[
\alpha = \min\!\left(1,\frac{\pi(\theta')q(\theta \mid \theta')}{\pi(\theta)q(\theta' \mid \theta)}\right),
\]
ensuring detailed balance and preservation of the target distribution. 

\medskip

MCMC thus transforms stochastic dependence into a computational tool, allowing Bayesian posterior distributions that lack closed form to be approximated numerically. 

% Section .....................................................................................................

\section{From Markov structure to modern language models}

The idea that the next state depends only on a limited context has deep connections to models of sequential data, including early probabilistic language models based on Markov chains, where the probability of a word depended only on a fixed number of preceding words. 

\medskip

In such \(n\)-gram models,
\[
\mathbb{P}(w_t \mid w_{t-1},\dots,w_1)
\approx
\mathbb{P}(w_t \mid w_{t-1},\dots,w_{t-n}),
\]
which is a finite-order Markov approximation. 

\medskip

Modern transformer-based large language models extend this idea by relaxing strict Markov assumptions and using attention mechanisms to condition on the entire preceding sequence, yet they remain fundamentally probabilistic models that assign conditional probabilities to tokens given context. 

\medskip

Thus, while transformers are not Markov chains in the classical sense, they inherit the foundational probabilistic principle that complex sequences can be modeled through structured conditional dependence, illustrating the enduring relevance of stochastic process theory in contemporary machine learning. 

% % Exercises ............................................................................................................

% \newpage

% \subsection*{Exercises}

% \textbf{1.} Exercise [...].\\

% \textbf{2.} Exercise [...].\\

% \textbf{3.} Exercise [...].\\
